\documentclass[10pt]{extarticle}
\usepackage[a4paper,margin=1.5cm]{geometry}
\usepackage[T1]{fontenc}
\usepackage{amsmath,amssymb}
\usepackage{graphicx}
\usepackage[export]{adjustbox}
\usepackage{multicol}
\usepackage{enumitem}
\usepackage{booktabs,array,tabularx}
\setlength{\columnsep}{0.7cm}
\setlength{\parindent}{0pt}
\setlength{\parskip}{3pt}
\pagestyle{empty}
\makeatletter
\newcommand{\qfigdef}[4]{\@namedef{qf@file@#1}{#2}\@namedef{qf@w@#1}{#3}\@namedef{qf@wide@#1}{#4}}
\newcommand{\qf@check}[1]{\@ifundefined{qf@file@#1}{\PackageError{score2}{Unknown figure id #1}{}}{}}
\newcommand{\qf@img}[4][]{\includegraphics[width=#2\dimexpr\@nameuse{qf@w@#4}pt\relax,max width=#3,max height=0.45\textheight,#1]{\@nameuse{qf@file@#4}}}
\newif\ifqf@full
\newif\ifqf@boxed
\newsavebox{\qf@box}
\newcommand{\qf@open}{\global\qf@fulltrue\par\global\qf@boxedtrue
  \setbox\qf@box\vbox\bgroup\hsize\textwidth\linewidth\textwidth}
\newcommand{\qf@emit}[1]{\par\egroup\global\qf@boxedfalse
  \ifdim\dimexpr\ht\qf@box+\dp\qf@box\relax<#1\textheight\noindent\usebox{\qf@box}\par\else\unvbox\qf@box\fi}
\newcommand{\qfullsplit}{\ifqf@boxed\qf@emit{0.85}\fi}
\newcommand{\qf@close}{\ifqf@boxed\qf@emit{0.45}\fi\global\qf@fullfalse}
\newcommand{\qfullbegin}{\end{multicols}\qf@open}
\newcommand{\qfullend}{\qf@close\begin{multicols}{2}}
\newcommand{\qfullstart}{\qf@open}
\newcommand{\qfulldone}{\qf@close}
\newcommand{\qfig}[2][1]{\qf@check{#2}%
  \ifqf@full\qf@blk{#1}{#2}\else
  \ifnum\@nameuse{qf@wide@#2}=1 \expandafter\@firstoftwo\else\expandafter\@secondoftwo\fi
  {\end{multicols}\begin{center}\qf@img{#1}{\textwidth}{#2}\end{center}\begin{multicols}{2}}%
  {\qf@blk{#1}{#2}}\fi}
\newcommand{\qf@blk}[2]{\par\smallskip{\centering\qf@img{#1}{\linewidth}{#2}\par}\smallskip}
\newcommand{\qfigi}[2][1]{\qf@check{#2}\qf@img[valign=c]{#1}{\linewidth}{#2}}
\makeatother
\newcommand{\papersec}[1]{\par\bigskip{\centering\Large\bfseries PAPER-#1\par}\medskip}
\newcommand{\partsec}[2]{\par\medskip\noindent{\large\bfseries PART-#1: #2}\par\smallskip}
\newcommand{\sect}[1]{\par\smallskip\noindent\textbf{#1}\par}
\newcommand{\qn}[1]{\par\smallskip\noindent\textbf{#1.}~}
\newcommand{\paraq}[2]{\par\smallskip\noindent\textbf{Paragraph for Questions #1 and #2}\par}
\newenvironment{qoptions}{\begin{itemize}[leftmargin=1.8em,itemsep=1pt,topsep=2pt,label={}]}{\end{itemize}}
\newcommand{\opt}[2]{\item[(#1)] #2}
% figure registry (generated by score2_tools.py finalize)
\qfigdef{p1_phy_s1_q01_1}{Score-I - Test 02 (31 Dec 2021)-img/p1_phy_s1_q01_1.png}{130.1}{0}
\qfigdef{p1_phy_s1_q02_1}{Score-I - Test 02 (31 Dec 2021)-img/p1_phy_s1_q02_1.png}{51.1}{0}
\qfigdef{p1_phy_s2_q05_1}{Score-I - Test 02 (31 Dec 2021)-img/p1_phy_s2_q05_1.png}{149.1}{0}
\qfigdef{p1_phy_s2_q08_1}{Score-I - Test 02 (31 Dec 2021)-img/p1_phy_s2_q08_1.png}{154.3}{0}
\qfigdef{p1_phy_s2_q09_1}{Score-I - Test 02 (31 Dec 2021)-img/p1_phy_s2_q09_1.png}{157.0}{0}
\qfigdef{p1_phy_s2_q10_1}{Score-I - Test 02 (31 Dec 2021)-img/p1_phy_s2_q10_1.png}{159.0}{0}
\qfigdef{p1_phy_s3_q01_1}{Score-I - Test 02 (31 Dec 2021)-img/p1_phy_s3_q01_1.png}{167.8}{0}
\qfigdef{p1_phy_s3_q05_1}{Score-I - Test 02 (31 Dec 2021)-img/p1_phy_s3_q05_1.png}{89.2}{0}
\qfigdef{p1_phy_s3_q06_1}{Score-I - Test 02 (31 Dec 2021)-img/p1_phy_s3_q06_1.png}{95.8}{0}
\qfigdef{p1_chem_s1_q01_A}{Score-I - Test 02 (31 Dec 2021)-img/p1_chem_s1_q01_A.png}{48.9}{0}
\qfigdef{p1_chem_s1_q01_B}{Score-I - Test 02 (31 Dec 2021)-img/p1_chem_s1_q01_B.png}{53.1}{0}
\qfigdef{p1_chem_s1_q01_C}{Score-I - Test 02 (31 Dec 2021)-img/p1_chem_s1_q01_C.png}{53.6}{0}
\qfigdef{p1_chem_s1_q01_D}{Score-I - Test 02 (31 Dec 2021)-img/p1_chem_s1_q01_D.png}{57.8}{0}
\qfigdef{p1_chem_s1_q01_5}{Score-I - Test 02 (31 Dec 2021)-img/p1_chem_s1_q01_5.png}{102.6}{0}
\qfigdef{p1_chem_s2_q08_C}{Score-I - Test 02 (31 Dec 2021)-img/p1_chem_s2_q08_C.png}{82.3}{0}
\qfigdef{p1_chem_s2_q09_1}{Score-I - Test 02 (31 Dec 2021)-img/p1_chem_s2_q09_1.png}{94.8}{0}
\qfigdef{p1_chem_s2_q09_A}{Score-I - Test 02 (31 Dec 2021)-img/p1_chem_s2_q09_A.png}{28.2}{0}
\qfigdef{p1_chem_s2_q09_B}{Score-I - Test 02 (31 Dec 2021)-img/p1_chem_s2_q09_B.png}{26.9}{0}
\qfigdef{p1_chem_s2_q09_C}{Score-I - Test 02 (31 Dec 2021)-img/p1_chem_s2_q09_C.png}{31.3}{0}
\qfigdef{p1_chem_s2_q09_D}{Score-I - Test 02 (31 Dec 2021)-img/p1_chem_s2_q09_D.png}{73.2}{0}
\qfigdef{p1_chem_s2_q10_A}{Score-I - Test 02 (31 Dec 2021)-img/p1_chem_s2_q10_A.png}{76.4}{0}
\qfigdef{p1_chem_s2_q10_B}{Score-I - Test 02 (31 Dec 2021)-img/p1_chem_s2_q10_B.png}{71.5}{0}
\qfigdef{p1_chem_s2_q10_C}{Score-I - Test 02 (31 Dec 2021)-img/p1_chem_s2_q10_C.png}{68.1}{0}
\qfigdef{p1_chem_s2_q10_D}{Score-I - Test 02 (31 Dec 2021)-img/p1_chem_s2_q10_D.png}{69.2}{0}
\qfigdef{p2_phy_s1_q07_1}{Score-I - Test 02 (31 Dec 2021)-img/p2_phy_s1_q07_1.png}{115.3}{0}
\qfigdef{p2_phy_s1_q09_1}{Score-I - Test 02 (31 Dec 2021)-img/p2_phy_s1_q09_1.png}{152.1}{0}
\qfigdef{p2_phy_s1_q13_1}{Score-I - Test 02 (31 Dec 2021)-img/p2_phy_s1_q13_1.png}{149.9}{0}
\qfigdef{p2_phy_s1_q15_1}{Score-I - Test 02 (31 Dec 2021)-img/p2_phy_s1_q15_1.png}{148.2}{0}
\qfigdef{p2_phy_s1_q15_A}{Score-I - Test 02 (31 Dec 2021)-img/p2_phy_s1_q15_A.png}{105.1}{0}
\qfigdef{p2_phy_s1_q15_B}{Score-I - Test 02 (31 Dec 2021)-img/p2_phy_s1_q15_B.png}{103.6}{0}
\qfigdef{p2_phy_s1_q15_C}{Score-I - Test 02 (31 Dec 2021)-img/p2_phy_s1_q15_C.png}{104.9}{0}
\qfigdef{p2_phy_s1_q15_D}{Score-I - Test 02 (31 Dec 2021)-img/p2_phy_s1_q15_D.png}{105.6}{0}
\qfigdef{p2_phy_s1_q16_1}{Score-I - Test 02 (31 Dec 2021)-img/p2_phy_s1_q16_1.png}{147.2}{0}
\qfigdef{p2_phy_s1_q17_1}{Score-I - Test 02 (31 Dec 2021)-img/p2_phy_s1_q17_1.png}{109.5}{0}
\qfigdef{p2_phy_s2_q04_1}{Score-I - Test 02 (31 Dec 2021)-img/p2_phy_s2_q04_1.png}{148.0}{0}
\qfigdef{p2_chem_s1_q15_1}{Score-I - Test 02 (31 Dec 2021)-img/p2_chem_s1_q15_1.png}{142.0}{0}
\qfigdef{p2_chem_s1_q15_A}{Score-I - Test 02 (31 Dec 2021)-img/p2_chem_s1_q15_A.png}{134.7}{0}
\qfigdef{p2_chem_s1_q15_B}{Score-I - Test 02 (31 Dec 2021)-img/p2_chem_s1_q15_B.png}{131.5}{0}
\qfigdef{p2_chem_s1_q15_C}{Score-I - Test 02 (31 Dec 2021)-img/p2_chem_s1_q15_C.png}{133.0}{0}
\qfigdef{p2_chem_s1_q15_D}{Score-I - Test 02 (31 Dec 2021)-img/p2_chem_s1_q15_D.png}{135.6}{0}
\qfigdef{p2_chem_s1_q17_A}{Score-I - Test 02 (31 Dec 2021)-img/p2_chem_s1_q17_A.png}{33.5}{0}
\qfigdef{p2_chem_s1_q17_B}{Score-I - Test 02 (31 Dec 2021)-img/p2_chem_s1_q17_B.png}{33.5}{0}
\qfigdef{p2_chem_s1_q17_C}{Score-I - Test 02 (31 Dec 2021)-img/p2_chem_s1_q17_C.png}{33.5}{0}
\qfigdef{p2_chem_s1_q17_D}{Score-I - Test 02 (31 Dec 2021)-img/p2_chem_s1_q17_D.png}{33.7}{0}
\qfigdef{p2_chem_s1_q18_1}{Score-I - Test 02 (31 Dec 2021)-img/p2_chem_s1_q18_1.png}{145.2}{0}
\qfigdef{p2_chem_s1_q19_1}{Score-I - Test 02 (31 Dec 2021)-img/p2_chem_s1_q19_1.png}{151.1}{0}
\qfigdef{p2_chem_s1_q20_A}{Score-I - Test 02 (31 Dec 2021)-img/p2_chem_s1_q20_A.png}{128.6}{0}
\qfigdef{p2_chem_s1_q20_B}{Score-I - Test 02 (31 Dec 2021)-img/p2_chem_s1_q20_B.png}{128.6}{0}
\qfigdef{p2_chem_s1_q20_C}{Score-I - Test 02 (31 Dec 2021)-img/p2_chem_s1_q20_C.png}{130.6}{0}
\qfigdef{p2_chem_s1_q20_D}{Score-I - Test 02 (31 Dec 2021)-img/p2_chem_s1_q20_D.png}{134.7}{0}
\qfigdef{p2_chem_s2_q01_1}{Score-I - Test 02 (31 Dec 2021)-img/p2_chem_s2_q01_1.png}{147.5}{0}
\qfigdef{p2_chem_s2_q02_1}{Score-I - Test 02 (31 Dec 2021)-img/p2_chem_s2_q02_1.png}{121.8}{0}

\begin{document}
\begin{center}{\Large\bfseries JEE (Main+Advanced) Score-I --- Test 02 (31 Dec 2021)}\end{center}
\begin{multicols}{2}
\papersec{1}
\partsec{1}{PHYSICS}
\sect{Section-I(i)}
\qn{1} A mercury barometer has a cylindrical capillary tube of radius $r=\sqrt2$ mm inverted vertically on the mercury in the trough. Atmospheric pressure is 76 cm of mercury but the level of mercury in the capillary is 75cm. The portion of the tube without mercury is 10cm in length. Considering that this region is occupied by ideal gasses, the number of moles of the ideal gas will be :

[Given : Density of mercury = 13.6 g/cc, surface tension of mercury = 0.5 N/m, temperature = $37^\circ$C, $g=10\ \mathrm{m/s^2}$, universal gas constant = $R$, angle of contact of mercury on tube = $135^\circ$]
\qfig{p1_phy_s1_q01_1}
\begin{qoptions}
\opt{A}{$0.55\dfrac{\pi}{R}\times10^{-6}$}
\opt{B}{$1.2\dfrac{\pi}{R}\times10^{-6}$}
\opt{C}{$0.88\dfrac{\pi}{R}\times10^{-6}$}
\opt{D}{$1.6\dfrac{\pi}{R}\times10^{-6}$}
\end{qoptions}
\qn{2} A slender bar of weight $W_0$ and Young's modulus $Y$ and coefficient of thermal expansion $\alpha$ is suspended from upper end. It just touches ground without generating a contact force. Now temperature is raised by $\Delta T$. The value of $\Delta T$ so that there is compression in the whole bar is
\qfig{p1_phy_s1_q02_1}
\begin{qoptions}
\opt{A}{$\dfrac{W_0}{2YA\alpha}$}
\opt{B}{$\dfrac{2W_0}{YA\alpha}$}
\opt{C}{$\dfrac{W_0}{YA\alpha}$}
\opt{D}{$\dfrac{2W_0}{3YA\alpha}$}
\end{qoptions}
\qn{3} A thin metallic shell of negligible thickness and radius `a' has water at $0^\circ$C completely filled inside it. This metallic shell is kept in a environment of $-\theta^\circ$C. The time after which $\dfrac18^{\text{th}}$ of the volume of water is left and rest of the water freezes, will be : [Given : conductivity of ice $K$, Latent heat of fusion $L$, density of ice $\rho$)
\begin{qoptions}
\opt{A}{$\dfrac{\rho La^2}{12K\theta}$}
\opt{B}{$\dfrac{\rho La^2}{6K\theta}$}
\opt{C}{$\dfrac{\rho La^2}{4K\theta}$}
\opt{D}{$\dfrac{\rho La^2}{16K\theta}$}
\end{qoptions}
\qn{4} Two objects A and B are kept hanging in the same environment. Object A is a spherical shell of radius $R$ and object B is an annular disc of inner radius $R$ and outer radius $2R$. Width of the sheet of shell and width of disc are same. Density and specific heat of object A and B are respectively $\rho_A$, $S_A$ and $\rho_B$, $S_B$. If initially the temperatures of A \& B are same and the rate at which their temperatures come down are same, the correct relation will be :
\begin{qoptions}
\opt{A}{$\rho_AS_A=\rho_BS_B$}
\opt{B}{$\rho_AS_A=2\rho_BS_B$}
\opt{C}{$2\rho_AS_A=3\rho_BS_B$}
\opt{D}{$2\rho_AS_A=\rho_BS_B$}
\end{qoptions}
\sect{Section-I(ii)}
\qn{5} A student studied the photoelectric effect from two different metals by performing two separate experiments. In the first experiment using a light source where frequency of the light can be varied, the student obtained the following plot of stopping potential versus frequency. In the second experiment using light source of frequency $1.2f_0$, the student found the stopping potential for photoelectrons released from sample A is 2.4 V larger than that for sample B. Given : $f_0=1\times10^{15}$ Hz, Work function of metal of sample A is $\phi_A$ and that of sample B is $\phi_B$.
\qfig{p1_phy_s2_q05_1}
\begin{qoptions}
\opt{A}{$\phi_A=1.74$ eV}
\opt{B}{$C=0.42$}
\opt{C}{In second experiment $V_S$ for electrons from sample A is 3.23 volts.}
\opt{D}{In second experiment $V_S$ for electrons from sample B is 0.83 volts.}
\end{qoptions}
\qn{6} According to Bohr model, point like electrons orbit point like nucleus at particular, fixed radii. In a particular such atom with a nucleus of charge $Q=1.6\times10^{-18}$ C initially one electron is in orbital 1 at radial distance $r_1=47.7$ \AA\ from nucleus. The orbital 2 is at a radial distance $r_2=84.8$\AA. The electric force on electron in orbital 1 is $F_1$ and that in orbital 2 is $F_2$. The KE of electron in orbital 1 is $K_1$ and in orbital 2 is $K_2$.
\begin{qoptions}
\opt{A}{The ratio of $\dfrac{F_1}{F_2}$ is greater than 3}
\opt{B}{$9K_1=16K_2$}
\opt{C}{When electron does a transition from orbital 1 to 2 it absorbs energy equal to 0.66 eV}
\opt{D}{$K_1-K_2=0.66$ eV}
\end{qoptions}
\qn{7} A radioactive isotope with decay constant $\lambda$, starts leaking from a nuclear reactor at $t=0$ with rate $R(t)$, where $R(t)=R_0\,(1-e^{-Kt})$. Initially there were no atoms present outside the reactor (at $t=0$).
\begin{qoptions}
\opt{A}{The number of radioactive atoms present outside the reactor at any time $t$ is $n(t)=R_0\left[\dfrac1\lambda\left(1-e^{-\lambda t}\right)+\dfrac{1}{\lambda-K}\left(e^{-\lambda t}-e^{-Kt}\right)\right]$.}
\opt{B}{The number of radioactive atoms present outside the reactor after a long time is $n(t)=\dfrac{R_0}{\lambda}$.}
\opt{C}{The number of radioactive atoms present outside the reactor at any time $t$ is $n(t)=R_0\left[\dfrac{1}{\lambda-K}\left(1-e^{-\lambda t}\right)+\dfrac{1}{\lambda-K}\left(e^{-\lambda t}-e^{-Kt}\right)\right]$.}
\opt{D}{The number of radioactive atoms present outside the reactor after a long time is $n(t)=\dfrac{R_0}{\lambda-K}$.}
\end{qoptions}
\qn{8} A cyclic process A $\to$ B $\to$ C $\to$ A brings a diatomic ideal gas back to its initial state in a three step process shown. (take 1 atm = $10^5$Pa)
\qfig{p1_phy_s2_q08_1}
\begin{qoptions}
\opt{A}{Total work done by the gas in one full cycle is 150 KJ}
\opt{B}{Temperature of the gas keeps on increasing during process A $\to$ B}
\opt{C}{Temperature of the gas keeps on decreasing during process C $\to$ A}
\opt{D}{Heat rejected by the gas to the surrounding in process C $\to$ A is 1075 kJ}
\end{qoptions}
\qn{9} Water is being heated in an electrical heater. At a certain moment a piece of ice at temperature $0^\circ$C was put into the kettle. The figure depicts the dependence of water temperature as a function of time. The power of the kettle is $P=1$ kW. The latent heat of melting of ice $L=335$ KJ/ kg and specific heat capacity of water is 4.2 KJ/ kg K.
\qfig{p1_phy_s2_q09_1}
\begin{qoptions}
\opt{A}{Mass of water before putting the ice is 476 gm}
\opt{B}{Mass of ice is 19 gm}
\opt{C}{Temperature of the contents at $t=80$sec is $69.6^\circ$C}
\opt{D}{Mass of ice is 24gm.}
\end{qoptions}
\qn{10} Three cylindrical beams of same material are loaded as shown. The ratio of square of radii is $r_A^2:r_B^2:r_C^2=1:1.5:2$.
\qfig{p1_phy_s2_q10_1}
\begin{qoptions}
\opt{A}{The stress in the beams are in the ratio of 2 : 4 : 1.}
\opt{B}{The compression in the beams are in the ratio of 8 : 12 : 1.}
\opt{C}{The elastic energy per unit volume in the beams are in the ratio of 16 : 16 : 1.}
\opt{D}{The total elastic energy stored in the beams are in the ratio of 16 : 9 : 1}
\end{qoptions}
\qn{11} Consider a helium atom following Bohr model. Model the two electrons in the same circular orbit and as far as possible. Take `$m$' and `$e$' as the mass and charge of $e^-$. Then for the atom
\begin{qoptions}
\opt{A}{The ground state radius will be $\dfrac{4h^2\epsilon_0}{7\pi me^2}$}
\opt{B}{The first ionisation energy of this atom will be $\dfrac{41}{128}\dfrac{me^4}{h^2\epsilon_0^2}$}
\opt{C}{Taking the reference of potential energy at infinity as zero, the energy of an electron is negative of twice of kinetic energy.}
\opt{D}{The second ionisation energy of this atom will be $\dfrac{35}{256}\dfrac{me^4}{h^2\epsilon_0^2}$}
\end{qoptions}
\qn{12} A piston-cylinder contains 1 kg of ideal gas with $C_V=615$ J/Kgk and $R=410$ J/Kgk. The initial pressure of the gas is $P_1=200$ kPa. The piston-cylinder also contains a metal block (in compressible). The block has mass 2 kg and specific heat capacity $C_b=305$ J/kgk. The block and gas is initially at 300 K. Assume no heat transfer to/from the cylinder/piston. The piston is pushed until temperature of gas reaches 363 K. The process happens rapidly so that there is no heat transfer between the gas and the block of metal. The piston is then locked in place and the gas and block come to thermal equilibrium over time. (Take : $1.1^5=1.61$)
\begin{qoptions}
\opt{A}{The pressure of the gas just after the rapid compression is 322 kPa}
\opt{B}{The final temperature of the block is close to 332 K}
\opt{C}{The final pressure of the gas is 221 kPa}
\opt{D}{Work done to push the piston is proportional to the molar mass of the gas.}
\end{qoptions}
\sect{Section-II}
\qn{1} A plane wall has a constant temperature maintained on either side. On one side it is 500K and on the other side it is 300K.
\qfig{p1_phy_s3_q01_1}
Measurements show that a convex temperature distribution is present within the wall such that the mid-point temperature is 10 K higher than expected for a linear temperature distribution. Assuming thermal conductivity has a linear dependence on temperature, $K=K_0\,(1+\alpha T)$, the constant $\alpha$ has value $\dfrac1N\,\mathrm{K}^{-1}$, find the value $N$.
\qn{2} In a particular forensic investigation a body is found in an apartment where room temperature is $72^\circ$F. The body's temperature is taken at 9:00pm and is found to be $80^\circ$F and decreasing at a rate of $2^\circ$F per hour. Estimate the time (in hour) before which death occur. [Given : Temperature before dying = $98.6^\circ$F, $\ln(3.325)=1.20$]
\qn{3} A mixture of some $\mathrm{O_2}$ and 22 g $\mathrm{CO_2}$ is subjected to some change such that all the $\mathrm{O_2}$ molecules dissociated into atoms. The molar specific heat at constant volume decreases by $\dfrac23R$ after this process ($R$ is universal gas constant). The mass of $\mathrm{O_2}$ in the mixture in grams is $M$. Find $M$. (Take : Degree of freedom of $\mathrm{CO_2}$ as 5)
\qn{4} In a particular $\alpha$-decay of $\mathrm{U}^{238}_{92}$ nucleus, the amount of kinetic energy acquired by the $\alpha$ particle is $K$ MeV. Find $K$. [Given : Binding Energy of $\mathrm{U}^{238}_{92}$ = 7.5 MeV/nucleon., Binding Energy of $\mathrm{Th}^{234}_{90}$ = 7.6 MeV/nucleon, Binding Energy of $\mathrm{He}^{4}_{2}$ = 7.1 MeV/nucleon.]
\qn{5} In the shown joules Apparatus each block has a mass of 1.75 kg and the insulated tank is filled with 200 g of water. The block are allowed to fall through a distance of 3.50 m. If the energy lost in the bearings and through walls is neglected, find the rise in temperature of the water in degree celsius. [Specific heat capacity of water is 4200 J/kg$^\circ$c]
\qfig{p1_phy_s3_q05_1}
\qn{6} A punch for making holes in steel plates is shown in Figure. Assume that a punch having a diameter of 20 mm is used to punch a hole in a 8 mm thick plate, as shown in cross-sectional view. If a force $P=150$ KN is required, the average shear stress in a plate is close to $Q\times10^8\ \mathrm{N/m^2}$. Find the value of $Q$ ($Q$ is an integer)
\qfig{p1_phy_s3_q06_1}
\partsec{2}{CHEMISTRY}
\sect{Section-I(i)}
\qn{1} Poly(methyl acrylate) on alkaline hydrolysis gives
\begin{qoptions}
\opt{A}{\qfigi{p1_chem_s1_q01_A}}
\opt{B}{\qfigi{p1_chem_s1_q01_B}}
\opt{C}{\qfigi{p1_chem_s1_q01_C}}
\opt{D}{\qfigi{p1_chem_s1_q01_D}}
\end{qoptions}
\qn{2}
\qfig{p1_chem_s1_q01_5}
`X'. The electrophile involved in the reaction is
\begin{qoptions}
\opt{A}{$\mathrm{NO_2^-}$}
\opt{B}{$\mathrm{NO^+}$}
\opt{C}{$\mathrm{H_3O^+}$}
\opt{D}{$\mathrm{NO_2^+}$}
\end{qoptions}
\qn{3} Hybridisation of `P' in the product obtained from reaction of $\mathrm{PH_3}$ with HI is
\begin{qoptions}
\opt{A}{$sp^3$}
\opt{B}{$sp^3d$}
\opt{C}{$sp^3d^2$}
\opt{D}{Doesn't hybridise}
\end{qoptions}
\qn{4} Choose the only \textbf{INCORRECT} statement
\begin{qoptions}
\opt{A}{In stomach, the pepsin enzyme converts proteins into peptides while in intestine, the pancreatic trypsin converts proteins into amino acids by hydrolysis.}
\opt{B}{The catalytic reaction that depends upon the pore structure of the catalyst and the size of the reactant and product molecules is called shape-selective catalysis.}
\opt{C}{ZSM-5 converts alcohols directly into gasoline (petrol) by dehydrating them to give a mixture of hydrocarbons phase.}
\opt{D}{Heterogeneous catalysis may involve reactant and catalyst in the same phase but not necessarily products in same phase}
\end{qoptions}
\sect{Section-I(ii)}
\qn{5} Two solutions A and B form a negative deviation solution. The B.P. of pure A is less than pure B. Choose the \textbf{CORRECT} statements
\begin{qoptions}
\opt{A}{Boiling temperature of the mixture is always more than boiling point of A.}
\opt{B}{Boiling temperature of azeotropic mixture is always less than boiling point of B.}
\opt{C}{When a small amount of A is added to excess of B, B.P. of resulting mixutre increases.}
\opt{D}{Boiling point of azeotropic mixture is a fixed value at any external pressure}
\end{qoptions}
\qn{6} Choose \textbf{CORRECT} statement regarding properties of colloidal solutions.
\begin{qoptions}
\opt{A}{Ultramicroscope may be used to determine size of colloidal particles}
\opt{B}{A mixture of milk and water appears blue when observed by reflected light}
\opt{C}{Mixture in option (B), appears red when observed by transmitted light.}
\opt{D}{Brownian movement is responsible for the stability of sols}
\end{qoptions}
\qn{7} Which of the following do not exist ?
\begin{qoptions}
\opt{A}{$\mathrm{He_2}$}
\opt{B}{$\mathrm{HeH^+}$}
\opt{C}{$\mathrm{Be_2}$}
\opt{D}{$\mathrm{HeH^-}$}
\end{qoptions}
\qn{8} Select \textbf{CORRECT} order of dipole moment for given molecular species.
\begin{qoptions}
\opt{A}{$\mathrm{NF_3}<\mathrm{NH_3}$}
\opt{B}{$\mathrm{CHCl_3}<\mathrm{CH_2Cl_2}$}
\opt{C}{\qfigi{p1_chem_s2_q08_C}}
\opt{D}{$\mathrm{HCN}>\mathrm{CO}$}
\end{qoptions}
\qn{9} Mark \textbf{CORRECT} colours obtained on reaction of given compounds with
\qfig{p1_chem_s2_q09_1}
\begin{qoptions}
\opt{A}{\qfigi{p1_chem_s2_q09_A} $\Rightarrow$ orange}
\opt{B}{\qfigi{p1_chem_s2_q09_B} $\Rightarrow$ yellow}
\opt{C}{\qfigi{p1_chem_s2_q09_C} $\Rightarrow$ Butter yellow}
\opt{D}{\qfigi{p1_chem_s2_q09_D} $\Rightarrow$ Orange / red}
\end{qoptions}
\qn{10} Mark the \textbf{INCORRECT} structures of L-Arabinose
\begin{qoptions}
\opt{A}{\qfigi{p1_chem_s2_q10_A}}
\opt{B}{\qfigi{p1_chem_s2_q10_B}}
\opt{C}{\qfigi{p1_chem_s2_q10_C}}
\opt{D}{\qfigi{p1_chem_s2_q10_D}}
\end{qoptions}
\qn{11} In which of the following, observed freezing point depression of solution is more than expected?
\begin{qoptions}
\opt{A}{Acetic acid in benzene}
\opt{B}{Milk powder in water}
\opt{C}{conc. $\mathrm{HNO_3}$ in conc. $\mathrm{H_2SO_4}$}
\opt{D}{NaCl in water}
\end{qoptions}
\qn{12} Which among the following may be classified as solid sol ?
\begin{qoptions}
\opt{A}{Pumice stone}
\opt{B}{foam rubber}
\opt{C}{Gemstone}
\opt{D}{Whipped cream}
\end{qoptions}
\sect{Section-II}
\qn{1} A cylindrical molecule `X' has dimensions-radius = 0.5\AA\ and height = 4\AA. `X' molecules are adsorbed on a platinum surface of dimensions [20cm $\times$ 30cm] . If adsorption capacity of Pt- surface is [0.1 $\mu$g of X/cm$^2$] and molar mass of `X' is 60g, then maximum $\mu$-moles of `X' that can be adsorbed on given Pt- surface shall be

[Take : $\pi=\dfrac{22}{7}$, $N_A=6\times10^{23}$] (Note : 1 $\mu$ -- mole = $10^{-6}$ moles)
\qn{2} The total vapour pressure of 5 mol\% solution of $\mathrm{H_2S}$ in water at 300K is 40 torr. The vapour pressure of pure water is 18 torr at this temperature. Estimate the total vapour pressure (in torr) of 6 mol \% solution of $\mathrm{H_2S}$. [Apply Henry's law and Raoult's law where necessary. Assume $\mathrm{H_2S}$ does not ionise.]
\qn{3} How many of the amino acid has 2 chiral centre

(i) Serine \qquad (ii) Proline

(iii) Threonine \qquad (iv) Leucine

(v) Valine \qquad (vi) Tyrosine

(vii) Aspartic acid \qquad (viii) Lysine
\qn{4} The number of odd $e^-$ species present in the following set are :

NO, $\mathrm{KO_2}$, $\mathrm{OF_2}$, $\mathrm{N_2O_4}$, $\mathrm{ClO_3}$.
\qn{5} How many of the following contains a peroxy linkage ?

Hydrogen peroxide, $\mathrm{K_3CrO_8}$, chromic acid, permanganic acid, perxenic acid, pernitric acid, $\mathrm{H_3PO_5}$, marshall's acid, caro's acid, oleum.
\qn{6} Identify number of planar species which have oxygen atom lone pair participation in back bonding.

$\mathrm{H_3BO_3}$, $\mathrm{H_2O}$, $\mathrm{OF_2}$, $\mathrm{Cl_2O}$, CO, $\mathrm{(CH_3)_2O}$, $\mathrm{(SiH_3)_2O}$, $\mathrm{CH_3NCO}$
\partsec{3}{MATHEMATICS}
\sect{Section-I(i)}
\qn{1} If $A=(a_{ij})_{3\times3}$ and $a_{ij}>0\ \forall\ i,j\in\{1,2,3\}$ and $|A|\ne0$ and the number of zero element in adjoint of $A$ is $\lambda$, Then the maximum value of $\lambda$ is
\begin{qoptions}
\opt{A}{0}
\opt{B}{1}
\opt{C}{2}
\opt{D}{3}
\end{qoptions}
\qn{2} If $|z-3i|\le5$, then minimum value of $|z+2|$ is equal to
\begin{qoptions}
\opt{A}{0}
\opt{B}{2}
\opt{C}{$5-\sqrt{13}$}
\opt{D}{1}
\end{qoptions}
\qn{3} The expression $-a^2x^2+bx-c$ is positive for atleast one real $x$ is $(a,b,c\in R,\ a\ne0)$ -
\begin{qoptions}
\opt{A}{$b^2<4a^2c$}
\opt{B}{$b^2>4a^2c$}
\opt{C}{$b^2=8a^2c$}
\opt{D}{$4b^2<a^2c$}
\end{qoptions}
\qn{4} Let $a,b,c$ be the sum, product and sum of the reciprocals of 37 consecutive terms of G.P. (common ratio of G.P. is not 1). If $b^m=\left(\dfrac ac\right)^n$, where $m,n$ are natural numbers and $m\ \&\ n$ are co-prime, then the digit at units place in $(m^n+n^m)$ is
\begin{qoptions}
\opt{A}{1}
\opt{B}{2}
\opt{C}{7}
\opt{D}{None of these}
\end{qoptions}
\sect{Section-I(ii)}
\qn{5} Let $P$ and $Q$ be two square matrices of same order such that $PQ=P$, $QP=Q$. Where $P$ and $Q$ are both skew symmetric matrices then $\det(P+Q)$ can be
\begin{qoptions}
\opt{A}{0}
\opt{B}{1}
\opt{C}{$-1$}
\opt{D}{2}
\end{qoptions}
\qn{6} Let $A(z_1)$ be a point lies on the circle whose centre is at origin and radius is 4 unit, $z_2$ be the midpoint of $O$ and $z_1$, where $O$ is origin. Let $\omega=e^{\frac{i2\pi}{21}}$ and $\displaystyle\sum_{r=1}^{20}|z_1+\omega^rz_2|^2=K$ then the sum of digit of $K$ is divisible by...
\begin{qoptions}
\opt{A}{3}
\opt{B}{4}
\opt{C}{5}
\opt{D}{7}
\end{qoptions}
\qn{7} If $x,y,z$ are non zero real numbers such that
{\footnotesize\[\begin{gathered}a^2+b^2+c^2+\frac{1}{x^2y^2z^2}\Big((x^2y+xy^2)^2+(y^2z+zy^2)^2\\ +(xz^2+z^2x)^2\Big)=\frac{-2}{xyz}\left(ayz(y+z)+bxz(x+z)+cxy(x+y)\right)\end{gathered}\]}
where $a,b,c\in R-\{1\}$ and the value of $\left[\left|\dfrac{1+a}{1-a}+\dfrac{1+b}{1-b}+\dfrac{1+c}{1-c}\right|\right]$ ([.] greatest integer function) divides
\begin{qoptions}
\opt{A}{2}
\opt{B}{3}
\opt{C}{7}
\opt{D}{5}
\end{qoptions}
\qn{8} The equation $a_{12}x^{12}+a_{11}x^{11}+a_{10}x^{10}+\ldots+a_0=0$ has all positive real roots and $a_{12}=1$, $a_{11}=-4$ and $a_0=\dfrac{1}{3^{12}}$, then
\begin{qoptions}
\opt{A}{$a_1=\dfrac{-1}{2^9}$}
\opt{B}{$a_1=\dfrac{-4}{3^{10}}$}
\opt{C}{$a_2=\dfrac{22}{3^9}$}
\opt{D}{$a_3=\dfrac{-220}{3^9}$}
\end{qoptions}
\qn{9} If $a,b,c\in R$ and all are positive such that $\dfrac{K(abc)}{a+b+c}\le(a+b)^2+(a+b+4c)^2$. Then the largest integral value of $K$ is divisible by
\begin{qoptions}
\opt{A}{4}
\opt{B}{5}
\opt{C}{10}
\opt{D}{15}
\end{qoptions}
\qn{10} If $\alpha,\beta,\gamma\ne n\pi$ such that $e^{i(2\alpha-\beta-\gamma)}+e^{i(2\beta-\gamma-\alpha)}+e^{i(2\gamma-\alpha-\beta)}=-1$. Then $|e^{i\alpha}+e^{i\beta}+e^{i\gamma}|^2$ can be
\begin{qoptions}
\opt{A}{1}
\opt{B}{2}
\opt{C}{3}
\opt{D}{4}
\end{qoptions}
\qn{11} If {\small\[A=\begin{bmatrix}x+xy+xyz&0&0\\0&y+yz+xyz&0\\0&0&z+zx+xyz\end{bmatrix},\]} $|A|\ne0$ and $xyz=1$ and sum of maximum and minimum value of $\mathrm{Tr}(A^{-1})$ (where $\mathrm{Tr}(A)$ denotes sum of diagonal elements) is $\dfrac pq$, then which of following options is/are correct ($p$ and $q$ are coprime)
\begin{qoptions}
\opt{A}{$p+q=3$}
\opt{B}{$p+q=5$}
\opt{C}{$p-q=1$}
\opt{D}{$3p-2q=4$}
\end{qoptions}
\qn{12} Given the sequence of numbers $a_1,a_2,a_3\ldots a_{1003}$ which satisfy $\dfrac{a_1}{1+a_1}=\dfrac{a_2}{3+a_2}=\dfrac{a_3}{5+a_3}=\ldots=\dfrac{a_{1003}}{2005+a_{1003}}$, then
\begin{qoptions}
\opt{A}{Sequence of terms will be in A.P.}
\opt{B}{Sequence of terms will be in G.P.}
\opt{C}{If fourth term of sequence 14 then sum of first five terms will be 50.}
\opt{D}{If fifth term of sequence 27 then sum of first ten terms will be 300.}
\end{qoptions}
\sect{Section-II}
\qn{1} If $a,b,c$ be three distinct natural numbers. Then the minimum value of
\[\left|a+be^{i\frac{2\pi}{3}}+ce^{i\frac{4\pi}{3}}\right|^2+\left|a+be^{-i\frac{2\pi}{3}}+ce^{-i\frac{4\pi}{3}}\right|^2\ \text{ is}\]
\qn{2} If $a,b,c$ be the distinct real numbers and $x,y,z$ be any three variables such that the system of linear equations

$(3z-2y)a+(2x-z)b+(y-3x)\,c=0$

$(y-3x)a+(3z-2y)b+(2x-z)\,c=0$

$(2x-z)a+(y-3x)b+(3z-2y)\,c=0$

can be express as $AX=O$, $A$ is $3\times3$ square matrix and $X=\begin{bmatrix}x\\y\\z\end{bmatrix}$ then the value of $\det(A)+\det(2A)+\det(3A)$ is :
\qn{3} If $\Delta=\begin{vmatrix}1&3\cos\theta&1\\\sin\theta&1&3\cos\theta\\1&\sin\theta&1\end{vmatrix}$, then maximum value of $\Delta$ is-
\qn{4} If $a,b,c$ are all natural numbers and all roots are not equal of the equation

$a^2x^4-2abx^3+(2ac+b^2)\,x^2-2bcx+c^2=0$ lies in $(0,1)$ then minimum value of $\dfrac{abc}{10}$ is
\qn{5} The value of $\displaystyle\sum_{n=2}^{\infty}\dfrac{n^4+3n^2+10n+10}{2^n(n^4+4)}$ is
\qn{6} If $\alpha_1,\alpha_2,\alpha_3\ldots,\alpha_{11}$ are roots of the equation $x^{11}+x^{10}+\ldots+x+1=0$ and the area of the polygon whose vertices are $1,\alpha_1,\alpha_2,\alpha_3\ldots,\alpha_{11}$ is $\lambda$ square unit. Then the value of $\lambda$ is
\papersec{2}
\partsec{1}{PHYSICS}
\sect{Section-I}
\qn{1} The radius of the orbit of an electron in Hydrogen-like atom is $4.5a_0$, where $a_0$ is the Bohr radius. Its orbital angular momentum is $\dfrac{3h}{2\pi}$. It is given that $h$ is Planck's constant and $R$ is Rydberg constant. The possible wavelength, when the atom de--excites, is
\begin{qoptions}
\opt{A}{$\dfrac{9}{31R}$}
\opt{B}{$\dfrac{9}{16R}$}
\opt{C}{$\dfrac{9}{5R}$}
\opt{D}{$\dfrac{4}{3R}$}
\end{qoptions}
\qn{2} The frequency of light emitted for the transition $n=4$ to $n=2$ of $\mathrm{He^+}$ is equal to the transition in H atom corresponding to which of the following ?
\begin{qoptions}
\opt{A}{$n=3$ to $n=1$}
\opt{B}{$n=3$ to $n=2$}
\opt{C}{$n=4$ to $n=3$}
\opt{D}{$n=2$ to $n=1$}
\end{qoptions}
\qn{3} The nuclide I (mass number 131) is radioactive, with a half-life of 8 days. At noon on march 1 the activity of a certain sample is 600 Bq. The activity at noon on march 24(same year) will be
\begin{qoptions}
\opt{A}{Little less than 200 Bq}
\opt{B}{Little more than 200 Bq}
\opt{C}{Little less than 75 Bq}
\opt{D}{Little more than 75 Bq}
\end{qoptions}
\qn{4} A U(mass number 238) nucleus emits a 4.196 MeV alpha particle. Calculate the disintegration energy $Q$ for this process, taking the recoil energy of the residual $^{234}$Th nucleus into account.
\begin{qoptions}
\opt{A}{4.196 MeV}
\opt{B}{4.267 MeV}
\opt{C}{4.401 MeV}
\opt{D}{4.200 MeV}
\end{qoptions}
\qn{5} One mole of gas A, with molecular diameter $2d$ and average molecular speed $v$, is placed inside a certain container. Two moles of gas B, with molecular diameter $d$ and average molecular speed $2v$, is placed in an identical container. Ratio of average collision rate of gas A within its container to the average collision rate of gas B within its container is
\begin{qoptions}
\opt{A}{1 : 3}
\opt{B}{1 : 1}
\opt{C}{3 : 1}
\opt{D}{2 : 1}
\end{qoptions}
\qn{6} We know that for an adiabatic process $PV^\gamma=a$ constant. Evaluate ``a constant'' for an adiabatic process involving exactly 10.0 mol of an ideal gas passing through the state having exactly $P=1.0$ atm and $T=1200$ K. Assume a diatomic gas whose molecules rotate but do not oscillate. [1 atm = $10^5$pa]
\begin{qoptions}
\opt{A}{4000 SI units}
\opt{B}{10000 SI units}
\opt{C}{100000 SI units}
\opt{D}{40000 SI units}
\end{qoptions}
\qn{7} In figure, the change in internal energy of a gas that is taken from A to C is 800 J. The work done along the path ABC is 500 J. If the pressure at point A is five times that at point C, what is the work done by the system in going from C to D (path 1) ?
\qfig{p2_phy_s1_q07_1}
\begin{qoptions}
\opt{A}{+100 J}
\opt{B}{$-100$ J}
\opt{C}{+80 J}
\opt{D}{$-80$ J}
\end{qoptions}
\qn{8} In one particular experiment of flow calorimeter, a liquid with a density of 800 $\mathrm{Kg/m^3}$ flows through the calorimeter at the rate of 4.00 $\mathrm{cm^3}$/s. At steady state, a temperature difference of $0.5^\circ$C is established between the input and output points when energy is supplied by heat at the rate of 3.2 J/s. What is the specific heat of the liquid?
\begin{qoptions}
\opt{A}{2000 J/Kg$^\circ$C}
\opt{B}{2400 J/Kg$^\circ$C}
\opt{C}{200 J/Kg$^\circ$C}
\opt{D}{20 J/Kg$^\circ$C}
\end{qoptions}
\qn{9} Consider two rods A and B of equal cross-sectional area $A$ are joined end to end as shown in figure. Initially two rods under thermal equilibrium with surrounding temperature of $T$ K. Length of the rods A, B are $\ell$, $2\ell$ respectively and their thermal conductivities are $2k$, $k$ respectively. If the ends of the rods are kept in contact of heat sources maintained at temperature $2T$ K and $3T$ K, the temperature of the junction at any time before achieving the steady state can be (lateral surface area of rods properly insulated from the surroundings)
\qfig{p2_phy_s1_q09_1}
\begin{qoptions}
\opt{A}{0.8 T}
\opt{B}{1.6 T}
\opt{C}{2.4 T}
\opt{D}{Data is not sufficient}
\end{qoptions}
\qn{10} A solid cube of edge length $r$, a solid sphere of radius $r$, and a solid hemisphere of radius $r$, all made of the same material, are maintained at temperature 300 K, 350 K and 400 K respectively, in an environment at temperature 350 K. Rank the objects according to the magnitude of net rate at which thermal radiation is exchanged with the environment, greatest first.
\begin{qoptions}
\opt{A}{Cube, solid hemisphere, solid sphere}
\opt{B}{Solid sphere, solid hemisphere, cube}
\opt{C}{Solid hemisphere, solid sphere, cube}
\opt{D}{Solid hemisphere, cube, solid sphere}
\end{qoptions}
\qn{11} Two spherical bodies A (radius 6 cm) and B (radius 18 cm) are at temperature $T_1$ and $T_2$ respectively. The maximum intensity in the emission spectrum of A is at 500 nm and in that of B is at 1500 nm. Considering them to be black bodies, what will be the ratio of the rate of total energy radiated by A to that of B ?
\begin{qoptions}
\opt{A}{3 : 1}
\opt{B}{1 : 3}
\opt{C}{1 : 9}
\opt{D}{9 : 1}
\end{qoptions}
\qn{12} An aluminum cup of 100 $\mathrm{cm^3}$ capacity is completely filled with glycerin at $22^\circ$C. How much glycerin, if any, will spill out of the cup if the temperature of both the cup and the glycerin is increased to $33.5^\circ$C? (The coefficient of volume expansion of glycerin is $5.1\times10^{-4}/\mathrm{C^\circ}$ and coefficient of linear expansion of aluminium is $75\times10^{-6}/\mathrm{C^\circ}$
\begin{qoptions}
\opt{A}{0.26 $\mathrm{cm^3}$}
\opt{B}{0.50 $\mathrm{cm^3}$}
\opt{C}{0.20 $\mathrm{cm^3}$}
\opt{D}{0.30 $\mathrm{cm^3}$}
\end{qoptions}
\qn{13} Figure shows stress versus strain plot for aluminium wire that is stretched gradually by a machine pulling in opposite directions at the two ends of the wire. The scale of the stress axis is set by $S=7.0$ in units of $10^7\ \mathrm{N/m^2}$. The wire has an initial length of 0.800 m and an initial cross-sectional area of $2.00\times10^{-6}\,\mathrm{m^2}$. How much work does the force from the machine do on the wire to produce a strain of $1.00\times10^{-3}$ ?
\qfig{p2_phy_s1_q13_1}
\begin{qoptions}
\opt{A}{56 mJ}
\opt{B}{112 mJ}
\opt{C}{5000 J}
\opt{D}{500 J}
\end{qoptions}
\qn{14} Light emitting diode is constructed using p-n junction ($E_g=1.5$ eV). What is the wavelength of light emitted by this LED ?
\begin{qoptions}
\opt{A}{720 nm}
\opt{B}{560 nm}
\opt{C}{450 nm}
\opt{D}{827 nm}
\end{qoptions}
\qn{15} The logic circuit shown below is equivalent to:
\qfig{p2_phy_s1_q15_1}
\begin{qoptions}
\opt{A}{\qfigi{p2_phy_s1_q15_A}}
\opt{B}{\qfigi{p2_phy_s1_q15_B}}
\opt{C}{\qfigi{p2_phy_s1_q15_C}}
\opt{D}{\qfigi{p2_phy_s1_q15_D}}
\end{qoptions}
\qn{16} Find $I$ ? (knee voltage of silicon diode is 0.7 volt)
\qfig{p2_phy_s1_q16_1}
\begin{qoptions}
\opt{A}{2.24 mA}
\opt{B}{6.06 mA}
\opt{C}{0.21 mA}
\opt{D}{0.12 mA}
\end{qoptions}
\qn{17} Ideal gas used in heat engine undergoes a cyclic process, in which one process is isothermal and one processes is adiabatic. During the isothermal process 40 J heat is released by the gas and during the isochoric process 80 J heat is absorbed by the gas, find the efficiency of that heat engine
\qfig{p2_phy_s1_q17_1}
\begin{qoptions}
\opt{A}{0.8}
\opt{B}{0.3}
\opt{C}{0.4}
\opt{D}{0.5}
\end{qoptions}
\qn{18} In a real refrigerator, the low-temperature coils are at $-13^\circ$C, and the compressed gas in the condenser is at $26^\circ$C. What is the theoretical coefficient of performance?
\begin{qoptions}
\opt{A}{1}
\opt{B}{$\dfrac{20}{3}$}
\opt{C}{$\dfrac{3}{20}$}
\opt{D}{0.8}
\end{qoptions}
\qn{19} One mole of an ideal diatomic gas ($C_v=5$ calories /mol-k) was transformed from initial $25^\circ$C and 1L to the state when temperature is $100^\circ$C and volume 10L. The entropy change of the process can be expressed as ($R=2$ calories /mol / k) [ans in cal/k]:-
\begin{qoptions}
\opt{A}{$3\ell n\dfrac{298}{373}+2\ell n10$}
\opt{B}{$5\ell n\dfrac{373}{298}+2\ell n10$}
\opt{C}{$7\ell n\dfrac{298}{373}+2\ell n\dfrac{1}{10}$}
\opt{D}{$5\ell n\dfrac{373}{298}+2\ell n\dfrac{1}{10}$}
\end{qoptions}
\qn{20} 1 g ice at $-40^\circ$C is placed in a container having 1 g water at $20^\circ$C. Find equilibrium temperature (in $^\circ$C).

Latent heat of fusion of ice = 80 cal/gm

Specific heat capacity of water = 1 cal/gm$^\circ$C

Specific heat capacity of ice = 0.5 cal/g$^\circ$C
\begin{qoptions}
\opt{A}{$-20^\circ$C}
\opt{B}{$0^\circ$C}
\opt{C}{$15^\circ$C}
\opt{D}{$5^\circ$C}
\end{qoptions}
\sect{Section-II}
\qn{1} An insulated container of gas has two chambers separated by an insulating partition. One of the chambers has volume $V$ and contains ideal monoatomic gas at pressure $P$ and temperature $T$. The other chamber has volume $2V$ and contains diatomic ideal gas at pressure $3P$ and temperature $T$. If the partition is removed without doing any work on the gas, ratio of the final equilibrium temperature of the gas to the initial temperature of the monoatomic gas is :
\qn{2} Suppose the fractional efficiency of a cesium surface (with work function 1.80 eV) is $1.0\times10^{-16}$; that is, on average one electron is ejected for every $10^{16}$ photons that reach the surface. The current of electrons ejected from such a surface if it were illuminated with 620 nm light from a 2.00 mW laser and all the ejected electrons took part in the charge flow is $10^n$ A then $n$ is
\qn{3} A radioactive isotope, X undergoes simultaneous decay to different nuclei as:

X $\to$ P (half life 2 hr)

X $\to$ Q (half life 1 hr)

After 2 hours ratio of number of Q elements to number of X elements is (Assuming that initially neither P nor Q was present)
\qn{4} The figure shows the cross-section of the outer wall of a house built in a hill-resort to keep the house insulated from the freezing temperature of outside. The wall consists of teak wood of thickness $L_1$ and brick of thickness ($L_2=5L_1$), sandwitching two layers of an unknown material with identical thermal conductivities and thickness. The thermal conductivity of teak wood is $K$ and that of brick is ($K_2=5K$). Heat conduction through the wall has reached a steady state with the temperature of three surfaces being known. ($T_1=25^\circ$C, $T_2=20^\circ$C and $T_5=-20^\circ$C). Find $T_3+T_4$ (in $^\circ$C)
\qfig{p2_phy_s2_q04_1}
\qn{5} A steel railroad track has a length of 30.000 m when the temperature is $0.0^\circ$C, the ends of the rail are rigidly clamped at $0.0^\circ$C so that expansion is prevented. The thermal stress set up in the rail if its temperature is raised to $40.0^\circ$C is $n\times10^7\,\mathrm{N/m^2}$. Young's modules of steel is $20\times10^{10}\,\mathrm{N/m^2}$. Thermal linear coefficient of expansion of steel is $\left[11\times10^{-6}\,(^\circ\mathrm{C})^{-1}\right]$. Then $n$ is
\qn{6} A liquid takes 5 minutes to cool from $60^\circ$C to $50^\circ$C, when temperature of surrounding is $30^\circ$C. How long (in minute) will it take to cool from $50^\circ$C to $40^\circ$C? (Use average form of Newton's law of cooling)
\qn{7} Ratio of debroglie wavelength of electron in third orbit of hydrogen atom to the debroglie wavelength of electron in first orbit of hydrogen is
\qn{8} A cobalt target is bombarded with electrons, and the wavelengths of its characteristic x-ray spectrum are measured. There is also a second, fainter characteristic spectrum, which is due to an impurity in the cobalt. The wavelengths of the $\mathrm{K}_\alpha$ lines are 178.9 pm (cobalt) and 143.5 pm (impurity), and the proton number for cobalt is ($Z_{\mathrm{Co}}$) = 27. Determine the impurity proton number using only these data. $\left(\text{Take }\sqrt{\dfrac{178.9}{143.5}}=1.11\right)$
\qn{9} If voltage of AM signal varies from 9 mV to 14 mV, then percentage of modulation will be :
\qn{10} In amplitude modulation, modulation index is $\dfrac23$ and maximum amplitude of modulated wave is 10 then minimum amplitude of the modulated wave is
\partsec{2}{CHEMISTRY}
\sect{Section-I}
\qn{1} Consider both molecules ; $\mathrm{N(SiH_3)_3}$ and $\mathrm{NH(SiH_3)_2}$ have trigonal planar skeleton at nitrogen. Select the incorrect statement from the given below statements.
\begin{qoptions}
\opt{A}{$\mathrm{Si-\hat N-Si}$ bond angle in $\mathrm{NH(SiH_3)_2}$ $>$ $\mathrm{Si-\hat N-Si}$ bond angle in $\mathrm{N(SiH_3)_3}$}
\opt{B}{N--Si bond length in $\mathrm{NH(SiH_3)_2}$ $>$ N--Si bond length in $\mathrm{N(SiH_3)_3}$}
\opt{C}{N--Si bond length in $\mathrm{NH(SiH_3)_2}$ $<$ N--Si bond length in $\mathrm{N(SiH_3)_3}$}
\opt{D}{Basic strength in $\mathrm{NH(SiH_3)_2}$ $>$ basic strength in $\mathrm{N(SiH_3)_3}$}
\end{qoptions}
\qn{2} The dipole moment of NO is 0.16 D, and the distance between the atoms is 111 pm. The percent ionic character of NO bond is
\begin{qoptions}
\opt{A}{16 \%}
\opt{B}{8 \%}
\opt{C}{3 \%}
\opt{D}{5.38 \%}
\end{qoptions}
\qn{3} For $\mathrm{H_2O_2}$, $\mathrm{H_2S}$, $\mathrm{H_2O}$ and HF, the correct order of decreasing strength of hydrogen bonding is :
\begin{qoptions}
\opt{A}{$\mathrm{H_2O>HF>H_2O_2>H_2S}$}
\opt{B}{$\mathrm{HF>H_2O_2>H_2O>H_2S}$}
\opt{C}{$\mathrm{HF>H_2O>H_2O_2>H_2S}$}
\opt{D}{$\mathrm{H_2O_2>H_2O>HF>H_2S}$}
\end{qoptions}
\qn{4} Which of the following is the correct order of increasing covalent character -
\begin{qoptions}
\opt{A}{$\mathrm{NaCl<LiCl<BeCl_2}$}
\opt{B}{$\mathrm{BeCl_2<NaCl<LiCl}$}
\opt{C}{$\mathrm{BeCl_2<LiCl<NaCl}$}
\opt{D}{$\mathrm{LiCl<NaCl<BeCl_2}$}
\end{qoptions}
\qn{5} The common features of the species $\mathrm{N_2^{2-}}$, $\mathrm{O_2}$ and $\mathrm{NO^-}$ are :
\begin{qoptions}
\opt{A}{Bond order three and isoelectronic.}
\opt{B}{Bond order two and isoelectronic.}
\opt{C}{Bond order three but not isoelectronic.}
\opt{D}{Bond order two but not isoelectronic.}
\end{qoptions}
\qn{6} The correct order of increasing C -- O bond length of CO, $\mathrm{CO_3^{2-}}$, $\mathrm{CO_2}$ is-
\begin{qoptions}
\opt{A}{$\mathrm{CO_3^{2-}<CO_2<CO}$}
\opt{B}{$\mathrm{CO_2<CO_3^{2-}<CO}$}
\opt{C}{$\mathrm{CO<CO_3^{2-}<CO_2}$}
\opt{D}{$\mathrm{CO<CO_2<CO_3^{2-}}$}
\end{qoptions}
\qn{7} Which of the following molecules does not exist?
\begin{qoptions}
\opt{A}{$\mathrm{BH_3}$}
\opt{B}{$\mathrm{SnH_4}$}
\opt{C}{$\mathrm{SbH_5}$}
\opt{D}{$\mathrm{SiH_4}$}
\end{qoptions}
\qn{8} An ideal solution is formed by mixing two liquids `A' and `B'. Determine the pressure at which $X_A=Y_B$ ($P^\circ_A=400$ torr, $P^\circ_B=100$ torr).

($X_A$ and $Y_B$ are mole fractions of A and B in liquid and vapour forms, respectively, at equilibrium.)
\begin{qoptions}
\opt{A}{250 torr}
\opt{B}{200 torr}
\opt{C}{150 torr}
\opt{D}{125 torr}
\end{qoptions}
\qn{9} Neutral soaps which are more soluble in water than in oil, form which type of emulsion ?
\begin{qoptions}
\opt{A}{Oil in water}
\opt{B}{Water in oil}
\opt{C}{Water in water}
\opt{D}{None of the above}
\end{qoptions}
\qn{10} Acetic acid exists in benzene solution in the dimeric form. In an actual experiment the Van't Hoff factor was found to be 0.52. Calculate the degree of association and osmotic pressure of 0.1 M acetic acid solution at 300K
\begin{qoptions}
\opt{A}{0.96 and 9.6 atm}
\opt{B}{0.96 and 1.28 atm}
\opt{C}{0.46 and 4.83 atm}
\opt{D}{0.26 and 7.86 atm}
\end{qoptions}
\qn{11} Which among the following is correct statement -
\begin{qoptions}
\opt{A}{Spontaneous adsorption of gases on solid surface is an endothermic process as entropy decreases during adsorption}
\opt{B}{Formation of micelles takes place when temperature is below Kraft Temperature ($T_k$) and concentration is above critical micelle concentration (CMC)}
\opt{C}{A colloid of $\mathrm{Fe(OH)_3}$ is prepared by adding a little excess (required to completely precipitate $\mathrm{Fe^{3+}}$ ions as $\mathrm{Fe(OH)_3}$) of NaOH in $\mathrm{FeCl_3}$ solution. The particles of this sol will move towards anode during electrophoresis.}
\opt{D}{According to Hardy-Schulze rule for the coagulation of $\mathrm{As_2S_3}$ sol flocculation value of $\mathrm{Fe^{3+}}$ ion will be more than $\mathrm{Ba^{2+}}$ or $\mathrm{Na^+}$.}
\end{qoptions}
\qn{12} A binary solution of two liquid `A' \& `B' is formed. Liquid `B' is more volatile than liquid A. If external pressure over liquid mixture gradually decreases then
\begin{qoptions}
\opt{A}{Mole fraction of `A' in liquid phase increases}
\opt{B}{Mole fraction of `B' in liquid phase increases}
\opt{C}{Mole fraction of `A' in liquid phase decreases}
\opt{D}{Pressure fraction of `B' is increases}
\end{qoptions}
\qn{13} According to Henry's law, the partial pressure of gas ($P_g$) is proportional to mole fraction of gas in dissolved state i.e. $P_{\mathrm{gas}}=K_H\cdot X_{\mathrm{gas}}$ where $K_H$ is Henry's constant. Then which of the following statements is incorrect ?
\begin{qoptions}
\opt{A}{$K_H$ is characteristic constant for a given gas - solvent system at a given temperature.}
\opt{B}{Higher is the value of $K_H$, more is the solubility of gas for a given partial pressure of gas.}
\opt{C}{$K_H$ has temperature dependence}
\opt{D}{$K_H$ increases with increase in temperature}
\end{qoptions}
\qn{14} Promoters and Poison are the substances used in chemical reaction which
\begin{qoptions}
\opt{A}{Enhance and decrease the activity of catalyst respectively}
\opt{B}{Decrease and enhance the activity of catalyst respectively}
\opt{C}{Do not have any effect on catalyst}
\opt{D}{Can be used in place of catalyst whenever required}
\end{qoptions}
\qn{15} Identify the major product ?
\qfig{p2_chem_s1_q15_1}
\begin{qoptions}
\opt{A}{\qfigi{p2_chem_s1_q15_A}}
\opt{B}{\qfigi{p2_chem_s1_q15_B}}
\opt{C}{\qfigi{p2_chem_s1_q15_C}}
\opt{D}{\qfigi{p2_chem_s1_q15_D}}
\end{qoptions}
\qn{16} Benzene diazonium chloride when reacts with hypophosphorus acid produces which of the following compound ?
\begin{qoptions}
\opt{A}{Benzene}
\opt{B}{Phenyl phosphate}
\opt{C}{Phenol}
\opt{D}{Phenyl isocyanide}
\end{qoptions}
\qn{17} Which is most reactive towards nucleophilic aromatic substitution
\begin{qoptions}
\opt{A}{\qfigi{p2_chem_s1_q17_A}}
\opt{B}{\qfigi{p2_chem_s1_q17_B}}
\opt{C}{\qfigi{p2_chem_s1_q17_C}}
\opt{D}{\qfigi{p2_chem_s1_q17_D}}
\end{qoptions}
\qn{18} The correct statement about the following sugar X and Y is :
\qfig{p2_chem_s1_q18_1}
\begin{qoptions}
\opt{A}{Both (X) \& (Y) are non-reducing sugar}
\opt{B}{(X) is non reducing sugar \& (Y) is a reducing sugar}
\opt{C}{(X) is reducing sugar \& (Y) is non reducing sugar}
\opt{D}{Both (X) \& (Y) are reducing sugar}
\end{qoptions}
\qn{19} Consider the following Aldoses/Ketoses. Which of them can give same osazone on reaction with $\mathrm{PhNHNH_2}$
\qfig{p2_chem_s1_q19_1}
\begin{qoptions}
\opt{A}{I, II only}
\opt{B}{I, II and III only}
\opt{C}{I, II and IV only}
\opt{D}{I, II, III and IV}
\end{qoptions}
\qn{20} A tripeptide is written as Alanine-Glycine-Alanine. The correct structure of this tripeptide is :
\begin{qoptions}
\opt{A}{\qfigi{p2_chem_s1_q20_A}}
\opt{B}{\qfigi{p2_chem_s1_q20_B}}
\opt{C}{\qfigi{p2_chem_s1_q20_C}}
\opt{D}{\qfigi{p2_chem_s1_q20_D}}
\end{qoptions}
\sect{Section-II}
\qn{1}
\qfig{p2_chem_s2_q01_1}
Aspartame is 100 times as sweet as sucrose and is used as a sugar substitute.

How many of the following is(are) correct statement(s) about aspartame

I -- It is an ester derivative of dipeptide.

II -- It has amide linkage.

III -- It is a tripeptide.

IV -- It has four functional groups.
\qn{2} In following reaction scheme -
\qfig{p2_chem_s2_q02_1}
Molecular mass of the major product P is-
\qn{3} In detection of phosphorous by fusing the organic compound with $\mathrm{Na_2O_2}$ and then extracting it with water, we add $\mathrm{HNO_3}$ (conc.) and ammonium molybdate to get a yellow precipitate of phosphorus containing compound. The number of oxygen atoms present in the formula of the compound.
\qn{4} Molar mass of the sodium compound formed during sodium fusion extract of the organic compound containing N and S is
\qn{5} The maximum no. of atoms which could be present in a plane at an instant in $\mathrm{H_2SO_4}$ are `X' and in $\mathrm{H_3PO_4}$ are `Y'. Then the value of $X-\dfrac Y2$ is
\qn{6} How many of the following statements is/are \textbf{CORRECT} ?

(1) Boiling point order : $\mathrm{NH_3>SbH_3>AsH_3>PH_3}$

(2) Bond angle order : $\mathrm{NH_3>PH_3>AsH_3>SbH_3}$

(3) In both crystalline $\mathrm{NaHCO_3}$ and $\mathrm{KHCO_3}$, $\mathrm{HCO_3^-}$ forms only dimeric anion through hydrogen bond.

(4) Amongst $\mathrm{B_2}$, $\mathrm{C_2}$, $\mathrm{N_2^-}$ and $\mathrm{O_2}$; $\mathrm{N_2^-}$ and $\mathrm{O_2}$ on further ionization (losing single electron) forms thermodynamically more stable species.
\qn{7} How many of the following paramagnetic species have fractional bond order in their structure at room temperature?

NO, $\mathrm{O_2^+}$, $\mathrm{O_3}$, $\mathrm{O_2^-}$, $\mathrm{C_2^{2-}}$, $\mathrm{NO^+}$, $\mathrm{N_2^+}$
\qn{8} 560 gm, 2m Urea solution (freezing point, $T_f$) when kept at temperature which is $5^\circ$C below its freezing point ($T_f$), then half of the solvent freezes out. Calculate the weight of solvent that (in g) freezes out if the same solution is placed at temperature $15^\circ$C below its freezing point ($T_f$)
\qn{9} Gold number of Haemoglobin is 0.030. Hence 100 ml of gold sol will require how much haemoglobin (in mg) so that gold is not coagulated by 10 ml of 10\% NaCl solution
\qn{10} The vapour pressure of water at room temperature is lowered by 5\% by dissolving a solute in it. Then the molality of solution will be
\partsec{3}{MATHEMATICS}
\sect{Section-I}
\qn{1} Let $|z_1|=1$, $|z_2|=2$, $|z_3|=3$ and $z_1+z_2+z_3=3+\sqrt5\,i$. Then $\mathrm{Re}\left(z_1\bar z_2+z_2\bar z_3+z_3\bar z_1\right)$ is equal to
\begin{qoptions}
\opt{A}{$1$}
\opt{B}{$-1$}
\opt{C}{$-\dfrac12$}
\opt{D}{$0$}
\end{qoptions}
\qn{2} Let $f(x)=\begin{bmatrix}1&x&0\\0&1&0\\0&0&1\end{bmatrix}$. Then $\left(f(x)\right)^n$ is equal to
\begin{qoptions}
\opt{A}{$f(x)$}
\opt{B}{$n f(x)$}
\opt{C}{$f(nx)$}
\opt{D}{$f(x^n)$}
\end{qoptions}
\qn{3} If $a,b,c$ are unequal natural numbers, then the value of the determinant
\[\begin{vmatrix}a^2+b^2&ab+bc&ac+ba\\ba+cb&b^2+c^2&bc+ca\\ca+ab&cb+ac&c^2+a^2\end{vmatrix}\]
is
\begin{qoptions}
\opt{A}{$a+b+c$}
\opt{B}{$a+b$}
\opt{C}{$b+c$}
\opt{D}{$0$}
\end{qoptions}
\qn{4} Let $a_1=5$ be the first term of an A.P. and $S_n=\sum\limits_{k=1}^{n}a_k$. If $\dfrac{S_m}{S_n}$ is independent of $n$ and $m=5n$, then the sum of all possible values of $a_2$ is
\begin{qoptions}
\opt{A}{$3$}
\opt{B}{$20$}
\opt{C}{$5$}
\opt{D}{$15$}
\end{qoptions}
\qn{5} The range of $k$ for which $k\cos^2x-k\cos x+1\ge0$ for all $x\in\mathbb R$ is
\begin{qoptions}
\opt{A}{$k>-\dfrac12$}
\opt{B}{$k>4$}
\opt{C}{$-\dfrac12\le k\le4$}
\opt{D}{$\dfrac12\le k\le5$}
\end{qoptions}
\qn{6} If $(A+B)(A-B)=A^2-B^2$, then $\left(A^2BA^{-1}B^{-1}\right)^3$ is equal to
\begin{qoptions}
\opt{A}{$B^2$}
\opt{B}{$A^2$}
\opt{C}{$A^3$}
\opt{D}{$A^3B^3$}
\end{qoptions}
\qn{7} If $|z|\le2$, then the maximum value of $|iz+6-8i|$ is
\begin{qoptions}
\opt{A}{$10$}
\opt{B}{$8$}
\opt{C}{$12$}
\opt{D}{$6$}
\end{qoptions}
\qn{8} If $\begin{vmatrix}a&5x&p\\b&10y&5\\c&15z&15\end{vmatrix}=125$, then $\begin{vmatrix}3a&3b&c\\x&2y&z\\p&5&5\end{vmatrix}$ is equal to
\begin{qoptions}
\opt{A}{$25$}
\opt{B}{$50$}
\opt{C}{$75$}
\opt{D}{$100$}
\end{qoptions}
\qn{9} Let $S_n=1+\dfrac12+\dfrac1{2^2}+\cdots+\dfrac1{2^{n-1}}$. The least value of $n$ for which $2-S_n<\dfrac1{100}$ is
\begin{qoptions}
\opt{A}{$7$}
\opt{B}{$9$}
\opt{C}{$8$}
\opt{D}{$6$}
\end{qoptions}
\qn{10} The least value of $n\in\mathbb N$ for which $(n-2)x^2+8x+n+4>0$ for all $x\in\mathbb R$ is
\begin{qoptions}
\opt{A}{$3$}
\opt{B}{$5$}
\opt{C}{$4$}
\opt{D}{None of these}
\end{qoptions}
\qn{11} If $\dfrac{a}{|z_2-z_3|}=\dfrac{b}{|z_3-z_1|}=\dfrac{c}{|z_1-z_2|}$, then $\dfrac{a^2}{z_2-z_3}+\dfrac{b^2}{z_3-z_1}+\dfrac{c^2}{z_1-z_2}$ is equal to
\begin{qoptions}
\opt{A}{$0$}
\opt{B}{$abc$}
\opt{C}{$3abc$}
\opt{D}{$a+b+c$}
\end{qoptions}
\qn{12} The system of equations $2x+ay+6z=8$, $x+2y+bz=5$, $x+y+3z=4$ has a unique solution if
\begin{qoptions}
\opt{A}{$a\ne2,\ b\ne3$}
\opt{B}{$a\ne3,\ b\ne2$}
\opt{C}{$a\ne2,\ b=3$}
\opt{D}{$a=2,\ b\ne3$}
\end{qoptions}
\qn{13} If $T_n$ is the $n^{\text{th}}$ term of an H.P. and $\dfrac{T_2}{T_6}=9$, then $\dfrac{T_{10}}{T_4}$ is equal to
\begin{qoptions}
\opt{A}{$\dfrac{17}{5}$}
\opt{B}{$\dfrac{5}{17}$}
\opt{C}{$\dfrac{7}{19}$}
\opt{D}{$\dfrac{19}{7}$}
\end{qoptions}
\qn{14} Let $a\ne0$ and let $\alpha<-3$, $\beta<2$ be the roots of $ax^2+bx+c=0$. Then which of the following is always true?
\begin{qoptions}
\opt{A}{$a\left(a+|b|+c\right)>0$}
\opt{B}{$a\left(a+|b|+c\right)<0$}
\opt{C}{$9a-3b+c>0$}
\opt{D}{$(9a-3b+c)(4a+2b+c)<0$}
\end{qoptions}
\qn{15} If $\alpha,\beta$ are the roots of $x^2-2x+4=0$, then $\alpha^n+\beta^n$ is equal to
\begin{qoptions}
\opt{A}{$2^n\cos\dfrac{n\pi}{3}$}
\opt{B}{$2^{n+1}\cos\dfrac{n\pi}{3}$}
\opt{C}{$2^n\sin\dfrac{n\pi}{3}$}
\opt{D}{$2^{n+1}\sin\dfrac{n\pi}{3}$}
\end{qoptions}
\qn{16} If the system $x+ky+3z=0$, $3x+ky-2z=0$, $2x+3y-4z=0$ has a non-trivial solution $(x,y,z)$, then the value of $k$ is
\begin{qoptions}
\opt{A}{$\dfrac{11}{14}$}
\opt{B}{$-\dfrac{33}{2}$}
\opt{C}{$\dfrac{33}{20}$}
\opt{D}{$\dfrac{33}{2}$}
\end{qoptions}
\qn{17} If $\sin\theta+2\sin\phi+3\sin\psi=0$ and $\cos\theta+2\cos\phi+3\cos\psi=0$, then the value of $\cos3\theta+8\cos3\phi+27\cos3\psi=$
\begin{qoptions}
\opt{A}{$\cos(3\theta+3\phi+3\psi)$}
\opt{B}{$18\cos(\theta+\phi+\psi)$}
\opt{C}{$6\cos(\theta+\phi+\psi)$}
\opt{D}{$36\cos(\theta+\phi+\psi)$}
\end{qoptions}
\qn{18} The sum of the real values that the function $f(x)=\dfrac{x^2-3x+2}{x^2+x-6}$ can not take is
\begin{qoptions}
\opt{A}{$2$}
\opt{B}{$-3$}
\opt{C}{$\dfrac65$}
\opt{D}{$\dfrac45$}
\end{qoptions}
\qn{19} Sum of series $2+\dfrac57+\dfrac8{49}+\dfrac{11}{343}+\cdots\infty$ is
\begin{qoptions}
\opt{A}{$\dfrac{35}{12}$}
\opt{B}{$\dfrac{49}{18}$}
\opt{C}{$\dfrac{35}{16}$}
\opt{D}{$\dfrac7{12}$}
\end{qoptions}
\qn{20} For a complex number $Z$, if $|Z-1+i|+|Z+i|=1$, then the range of the principal argument of $Z$ is (where principal $\arg(Z)\in(-\pi,\pi]$)
\begin{qoptions}
\opt{A}{$\left[-\dfrac\pi4,\dfrac\pi4\right]$}
\opt{B}{$\left[\dfrac\pi4,\dfrac\pi2\right]$}
\opt{C}{$\left[-\dfrac\pi2,-\dfrac\pi4\right]$}
\opt{D}{$\left[-\dfrac\pi2,\dfrac\pi2\right]$}
\end{qoptions}
\sect{Section-II}
\qn{1} Let matrix $A=\begin{bmatrix}1&3\\2&1\end{bmatrix}$ and $b,c\in\mathbb R$ are such that $A^2+bA=cI_2$, then value of $b+c$ is
\qn{2} If $\begin{vmatrix}1&x&x^2\\x&x^2&1\\x^2&1&x\end{vmatrix}=3$, then value of $\begin{vmatrix}x^3-1&0&x-x^4\\0&x-x^4&x^3-1\\x-x^4&x^3-1&0\end{vmatrix}$ is
\qn{3} Let $\alpha,\beta,\gamma$ and $\delta$ be the roots (real or non real) of equations $x^4-3x+1=0$. The value of $\alpha^3+\beta^3+\gamma^3+\delta^3$ is
\qn{4} Let $z_1,z_2,z_3,z_4$ be the vertices of a square in the Argand plane which lie on the curve $|z-1-i|=3|z+1+i|$. If the area of the square is $\lambda$, then the value of $\lambda$ is
\qn{5} If $a+2b+3c=6$ and the greatest value of $abc^2$ is $k$, then the value of $16k$ is
\qn{6} Let $f(x)=x^2+(1-2\lambda)x+(\lambda^2-\lambda-2)$ have one root positive and the other root negative. Then the number of integral values of $\lambda$ is
\qn{7} If $A^{-1}=\begin{bmatrix}2&1&3\\0&4&5\\1&2&0\end{bmatrix}$, then $\left|\mathrm{adj}(3A)\right|$ is equal to
\qn{8} If $|z+5|\le4$ and $z(1+i)+\bar z(1-i)\ge-10$, then the maximum value of $|z+1|^2$ is $\alpha+\beta\sqrt2$. Then the value of $\alpha+\beta$ is
\qn{9} Let $a_1,a_2,a_3,\ldots$ be in G.P. Then the value of $\begin{vmatrix}\log a_n&\log a_{n+1}&\log a_{n+2}\\\log a_{n+3}&\log a_{n+4}&\log a_{n+5}\\\log a_{n+6}&\log a_{n+7}&\log a_{n+8}\end{vmatrix}$ is
\qn{10} Let $\langle a_n\rangle$ be a sequence such that $a_1=0$, $a_2=1$ and $a_n=a_{n-1}-a_{n-2}$ for $n\ge3$. Then the value of $a_{2016}$ is
\end{multicols}
\end{document}
